\section{Related Work}
\label{sec:related}

\todo{Comparison matrix as Table~\ref{tab:related}: rows = Yuho, Catala,
lam4, LexScript, Akoma Ntoso, LDOC, LegalRuleML; columns = expressivity
(deontic types, prioritised defaults, penalty combinators, illustration blocks),
tooling (parser, LSP, MCP, transpilers), verification hookup (Z3, Alloy,
none), and target coverage (single-statute / corpus / synthetic).}

\subsection{Catala}
Catala~\cite{catala2021} pioneered prioritised default rules for legal
encoding. We adopt the priority semantics directly for our exception
DAG (\S\ref{subsec:exceptions}) but disagree on top-level shape: Catala's
top-level construct is a \emph{scope} (a callable function), whereas
Yuho's top-level construct is a \emph{statute} block, which mirrors the
drafting unit and carries fields (definitions, illustrations, case law)
that a function-call shape does not naturally accommodate. Scope composition
is on our roadmap but as an additional layer, not a replacement for the
statute block.

\subsection{lam4}
lam4~\cite{lam4} centres on deontic logic and named-norm references. The
\texttt{IS\_INFRINGED} predicate is the obvious target for our planned
G10 inter-section semantic hookup. lam4 targets contracts and regulations;
Yuho targets penal code. The two grammars are non-overlapping but
philosophically aligned.

\subsection{LexScript}
LexScript~\cite{lexscript} is the closest tooling-shaped peer: tree-sitter
grammar, statute-shaped top-level. We borrow its Tarjan-SCC reference-graph
analysis (planned, blocked on G10). LexScript has a thinner element model
(no burden / no causal links / no Catala-style exception priority).

\subsection{Akoma Ntoso and LegalRuleML}
\todo{Position as orthogonal: markup standards for archival exchange,
not executable encodings. We can emit Akoma Ntoso XML as a future
transpile target without competing with it.}

\subsection{LDOC}
LDOC's authoring-loop emphasis (Word + PDF) is the inspiration for our
DOCX transpile target. Yuho diverges by not making Word the source of
truth: the canonical encoding is the \texttt{.yh} file in version control;
DOCX is one of six output formats.
